\section{MPI Matching in One Page}
\label{sec:background}

This section covers the parts of MPI the paper needs~\citep{mpi50}, for readers who know
concurrency but not MPI.

\paragraph{Ranks and messages.} An MPI program runs as $p$ processes, called
\emph{ranks}, numbered $0,\dots,p-1$, each with private memory. A send names a destination
rank, a \emph{tag} (an integer label) and a \emph{communicator} (a group of ranks with its own
message space). A receive names a source rank, a tag and a communicator; the source or the tag
may be a wildcard, \anysrc or \texttt{MPI\_ANY\_TAG}. A receive $r$ \emph{accepts} a send $s$,
written $r \accepts s$, when they agree on the communicator, $s$ is addressed to $r$'s rank,
and the source and tag of $r$ are wildcards or equal to those of $s$. A completed receive's
\emph{status} reports the actual source, tag and element count. Programs often branch on it,
for example to reply to the sender.

\paragraph{Matching and non-overtaking.} Matching pairs each receive with at most one
send and each send with at most one receive. Two ordering rules, together called
\emph{non-overtaking}~\citep[§3.5]{mpi50}, fix the rest. If a rank sends two messages that
the same receive accepts, the earlier send matches first. If a rank posts two receives that
both accept a message, the earlier receive takes it. The only freedom left is the arrival
order of messages from \emph{different} senders, and a wildcard receive is where it becomes
observable.

\paragraph{Blocking, nonblocking and completion.} \mpicall{Send} and \mpicall{Recv}
block. \mpicall{Isend} and \mpicall{Irecv} return a request at once and complete later,
in a \mpicall{Wait}, a \mpicall{Test} or one of their variants; \mpicall{Waitany}
returns whichever of several requests completes first. A receive is \emph{posted} when it is
issued, and matched then or when a matching message arrives, so a nonblocking receive is
matched long before its wait. \mpicall{Probe} and \mpicall{Iprobe} report the envelope
of a message that a receive would take, without receiving it.

\paragraph{Buffering.} A standard-mode send may complete before its receive is posted,
if the library copies the message into a buffer (\emph{eager} protocol), or wait until the
receive is posted (\emph{rendezvous}). Implementations choose by message size; Open MPI~4.1
sends up to \BimplOMPI bytes eagerly on shared memory by default. We model this as a
threshold $\Bt$: a standard send of more than $\Bt$ bytes blocks until matched.
\mpicall{Ssend} always behaves as $\Bt=0$, \mpicall{Bsend} as $\Bt=\infty$, and a ready send
(\mpicall{Rsend}) is erroneous unless the receive is already posted. Two ranks that both send
before they receive deadlock exactly when their messages exceed $\Bt$.

\paragraph{Collectives.} Collective operations (\mpicall{Barrier},
\mpicall{Bcast}, \mpicall{Reduce}, \dots) are called by every rank of a communicator in
the same order. The standard requires no more synchronisation than their data flow needs: a
\mpicall{Bcast} root may return before the others enter, and a \mpicall{Reduce} root waits
for everyone. We model each collective by the synchronisation its \emph{shape} guarantees.
